\documentclass[11pt]{article}

\usepackage[margin=1.15in]{geometry}
\usepackage{amsmath,amssymb,amsthm}
\usepackage[colorlinks=true,linkcolor=blue,citecolor=blue,urlcolor=blue]{hyperref}

\newtheorem{theorem}{Theorem}
\newtheorem{lemma}{Lemma}
\theoremstyle{remark}
\newtheorem*{remark}{Remark}

\newcommand{\Frac}{\operatorname{Frac}}
\newcommand{\Ann}{\operatorname{Ann}}
\newcommand{\Hom}{\operatorname{Hom}}
\newcommand{\Ext}{\operatorname{Ext}}
\newcommand{\soc}{\operatorname{soc}}
\newcommand{\rank}{\operatorname{rank}}

\title{Log-concavity of codimension-three level Hilbert functions\\ of type two}
\author{Isaac Zhang}
\date{\today}

\begin{document}

\maketitle

\begin{abstract}
Let $R = k[x_1,x_2,x_3]$, where $k$ is any field. We prove that the Hilbert
function $h = (1, 3, h_2, \dots, h_e = 2)$ of every standard graded Artinian
level quotient of $R$ of type two is log-concave. The argument uses the
minimal resolution of the reindexed graded Matlis dual. Its last syzygy has
rank one and full support; this forces every proper subset of the middle
columns to be independent over $\Frac(R)$. A putative failure of
log-concavity then reduces to a rank-one family of relations, which
determines a cyclic Gorenstein submodule. Stanley's characteristic-free
theorem for codimension-three Gorenstein Hilbert functions closes the
remaining case.
\end{abstract}

\section*{Introduction}

Let $A$ be a standard graded Artinian level algebra of codimension $r$ and
type $t$. Iarrobino \cite{Iarrobino} asked for which pairs $(r,t)$ the
Hilbert function of every such $A$ is log-concave, and answered it
positively for $r = 2$ and any $t$, and for $(r,t) = (3,1)$, and negatively
for $(r,t) = (4,1)$. Zanello \cite{ZanelloLC} settled every remaining pair
negatively, with one exception: $(r,t) = (3,2)$, which \cite{ZanelloLC}
records as open in every characteristic. Theorem \ref{thm:main} below
resolves that case.

Zanello \cite{ZanelloLC} raises the same question for pure $O$-sequences,
equivalently for the Hilbert functions of \emph{monomial} Artinian level
algebras, where $(3,2)$ was likewise the principal open case. That monomial
case has been settled by a combinatorial argument \cite{AlphaProofNexus}. It
does not imply the result proved here: the Hilbert function of $R/I$ is
unchanged on passing to a generic initial ideal, but levelness and type are
not preserved, since taking initial ideals can raise the socle type. The
implication runs the other way, as every pure $O$-sequence of codimension
three and type two is the Hilbert function of a monomial level algebra of
codimension three and type two.

The proof below has been formalized in Lean 4 with Mathlib; the development
is available at \url{https://github.com/tigerisaac/zanello-lc}. Its endpoint
takes the level algebra and Lemma \ref{lem:stanley} as its only hypotheses,
is free of \texttt{sorry}, and depends only on Lean's three standard
foundational axioms.

\section*{Notation and the one imported theorem}

For $j \ge 0$ set
\[
  N_j \;=\; \dim_k R_j \;=\; \binom{j+2}{2},
\]
and put $N_j = 0$ for $j < 0$. We will use the following
characteristic-free form of Stanley's theorem.

\begin{lemma}[Stanley] \label{lem:stanley}
If $B$ is a standard graded Artinian Gorenstein algebra of embedding
dimension at most three and socle degree $E$, then its Hilbert function is
symmetric and is nondecreasing through degree $\lfloor E/2 \rfloor$.
\end{lemma}

More precisely, the Hilbert function is an SI-sequence \cite{Stanley}; the
monotonicity stated above is the only part needed here. The result holds
over an arbitrary field. A characteristic-free proof is given in
\cite{ZanelloStanley}.

\section*{The theorem}

\begin{theorem} \label{thm:main}
Let $A = R/I$ be a standard graded Artinian level algebra of embedding
dimension three, socle degree $e$, and type two. Then
\[
  h_i^2 \;\ge\; h_{i-1} h_{i+1} \qquad (1 \le i \le e-1).
\]
The statement holds in every characteristic.
\end{theorem}

\begin{proof}
Put $s = e + 3$.

\medskip
\noindent\textbf{The dual complex.}
Since $A$ is Artinian, $\operatorname{depth} A = 0$, so the graded
Auslander--Buchsbaum formula gives $\operatorname{pd}_R A = 3$ and the
minimal free resolution of $A$ has the form
\[
  0 \longrightarrow G_3 \longrightarrow G_2 \longrightarrow G_1
    \longrightarrow R \longrightarrow A \longrightarrow 0 .
\]
Computing $\operatorname{Tor}^R_3(A,k)$ from the Koszul resolution of $k$ on
$x_1,x_2,x_3$ gives
\[
  \operatorname{Tor}^R_3(A,k) \;\cong\; \soc(A)(-3),
\]
while computing it from that resolution gives $G_3 \otimes_R k$. As $A$ is level
of type two with socle degree $e$, its socle is two-dimensional and
concentrated in degree $e$, so
\[
  G_3 \;=\; R(-e-3)^2 \;=\; R(-s)^2 .
\]

Reindex the graded Matlis dual by setting
\[
  M_d \;=\; \Hom_k(A_{e-d}, k), \qquad
  g_d \;=\; \dim_k M_d \;=\; h_{e-d} .
\]
By graded local duality $M \cong \Ext^3_R(A, R(-s))$, so applying
$\Hom_R(-, R(-s))$ to that resolution yields a minimal exact complex
\begin{equation} \label{eq:one}
  0 \longrightarrow R(-s) \xrightarrow{\;\delta_3\;} F_2
    \xrightarrow{\;\delta_2\;} F_1 \xrightarrow{\;\delta_1\;} R^2
    \longrightarrow M \longrightarrow 0,
\end{equation}
where $F_i = \Hom_R(G_{3-i}, R(-s))$ and the two outer terms are
$\Hom_R(R, R(-s)) = R(-s)$ and $\Hom_R(R(-s)^2, R(-s)) = R^2$.

Thus $M$ is generated by two elements in degree zero. Its socle is
one-dimensional and concentrated in degree $e$. Indeed, for $d < e$,
standard generation gives $A_{e-d} = R_1 A_{e-d-1}$, so no nonzero
functional in $M_d$ is annihilated by all of $R_1$.

Write
\[
  F_1 \;=\; \bigoplus_j R(-j)^{p_j}, \qquad
  F_2 \;=\; \bigoplus_j R(-j)^{q_j} .
\]

\medskip
\noindent\textbf{The circuit.}
Let $K = \Frac(R)$. Tensoring \eqref{eq:one} with $K$ shows that the kernel
of $\delta_2$ is one-dimensional. Its spanning vector has no zero
coordinate. Indeed, $\delta_3 = \Hom_R(\varepsilon, R(-s))$ where
$\varepsilon \colon G_1 \to R$ is the first differential of the resolution
of $A$, whose entries are a minimal homogeneous generating set of
$I$; dualizing, the image of a generator of $R(-s)$ under $\delta_3$ has
those same entries as its coordinates, and a minimal generator of a nonzero
ideal is nonzero. Consequently,
\begin{equation} \label{eq:two}
  \text{every proper subset of the columns of } \delta_2
  \text{ is } K\text{-linearly independent.}
\end{equation}

\medskip
\noindent\textbf{The rank estimate.}
Fix $d$ with $2 \le d \le e$. The corresponding log-concavity inequality is
the one centered at $g_{d-1}$. Define
\[
  Q_d \;=\; \sum_{j \le d} q_j, \qquad P_{<d} \;=\; \sum_{j < d} p_j ,
\]
and let $r_d$ be the $K$-rank of the restriction of $\delta_1$ to the
summands of $F_1$ having shift strictly less than $d$. Thus
$r_d \in \{0,1,2\}$. Finally, set
\[
  \varepsilon_d \;=\;
  \begin{cases}
    1, & Q_d = \rank F_2, \\
    0, & Q_d < \rank F_2 .
  \end{cases}
\]
By \eqref{eq:two}, the columns of $\delta_2$ having shift at most $d$ have
rank $Q_d - \varepsilon_d$. Since every differential in a minimal graded
free resolution has entries in the homogeneous maximal ideal, these columns
can have nonzero entries only in rows of $F_1$ whose shifts are strictly
less than $d$. Their images lie in the kernel of a rank-$r_d$ map from a
$P_{<d}$-dimensional $K$-space to $K^2$. Therefore
\begin{equation} \label{eq:three}
  Q_d - \varepsilon_d \;\le\; P_{<d} - r_d .
\end{equation}

Let $\Delta^2 g_d = g_d - 2g_{d-1} + g_{d-2}$. The Hilbert numerator of
\eqref{eq:one} is
\[
  (1-t)^3 \sum_{j \ge 0} g_j t^j
  \;=\; 2 - \sum_j p_j t^j + \sum_j q_j t^j - t^s .
\]
Since $d \le e < s$, summing its coefficients through degree $d$ gives
\[
  \Delta^2 g_d \;=\; 2 + Q_d - P_{<d} - p_d .
\]
Together with \eqref{eq:three}, this yields the central estimate
\begin{equation} \label{eq:four}
  \Delta^2 g_d \;\le\; 2 - r_d + \varepsilon_d - p_d .
\end{equation}
If $\Delta^2 g_d \le 0$, then $g_{d-2} + g_d \le 2 g_{d-1}$, and the
arithmetic--geometric mean inequality gives $g_{d-2} g_d \le g_{d-1}^2$.
Thus a failure of log-concavity would imply
\begin{equation} \label{eq:five}
  \Delta^2 g_d \;\ge\; 1 .
\end{equation}

\medskip
\noindent\textbf{Case analysis.}
First suppose $\varepsilon_d = 1$. Every summand $R(-b)$ of $F_2$ then has
$b \le d$. By the degree correspondence in the dual resolution, it comes
from a minimal generator of $I$ of degree $s - b \ge s - d$. Thus $I$ has no
nonzero element in degrees below $s-d$. With $m = e-d+2 = s-d-1$, the tested
triple is
\[
  (g_{d-2}, g_{d-1}, g_d) \;=\; (N_m, N_{m-1}, N_{m-2}).
\]
Since the ratios $N_j / N_{j-1}$ are nonincreasing, this triple is
log-concave.

We may therefore assume $\varepsilon_d = 0$. If $r_d = 2$, \eqref{eq:four}
gives $\Delta^2 g_d \le 0$, contradicting \eqref{eq:five}.

If $r_d = 0$, there are no summands of $F_1$ of shift below $d$. Indeed,
each such summand gives a nonzero minimal relation, so the restricted map
would have positive rank otherwise. Equation \eqref{eq:three} then gives
$Q_d = 0$. Consequently
\[
  g_{d-2} = 2N_{d-2} = d(d-1), \qquad g_{d-1} = 2N_{d-1} = d(d+1),
\]
and
\[
  g_d = 2N_d - p_d = (d+1)(d+2) - p_d .
\]
A direct calculation gives
\[
  g_{d-1}^2 - g_{d-2} g_d \;=\; d \bigl( 2(d+1) + (d-1) p_d \bigr) \;>\; 0 .
\]

\medskip
\noindent\textbf{The critical case $r_d = 1$.}
From \eqref{eq:four} and \eqref{eq:five}, a failure can occur only if
\begin{equation} \label{eq:six}
  p_d = 0, \qquad \Delta^2 g_d = 1 .
\end{equation}
All columns of $\delta_1$ having shift below $d$ span one line over $K$.
Because $R$ is a UFD, this line has a primitive homogeneous generator
$v = (v_1, v_2) \in R^2$, whose two components have a common degree $a < d$.
The submodule generated by the low-shift columns is then $vJ$ for a
homogeneous ideal $J \subseteq R$. To see this, divide one nonzero column by
the gcd of its two entries. The resulting primitive vector spans the line,
and Euclid's lemma shows that every other column is an $R$-multiple of it.

Because $p_d = 0$, no relation has shift $d$, while relations of shift
greater than $d$ have no component in degrees at most $d$. Therefore
\begin{equation} \label{eq:seven}
  \bigl( \ker(R^2 \to M) \bigr)_t \;=\; (vJ)_t \qquad (t \le d).
\end{equation}

Let $H$ be the image of $v$ in $M_a$ and put $n = d - a$.

If $H = 0$, then $v$ itself is a relation of degree $a < d$. Equation
\eqref{eq:seven} and the primitivity of $v$ force $1 \in J$, so $J = R$. In
this case set $B_j = 0$ for every $j$.

Suppose $H \ne 0$. The finite-length module $RH$ has nonzero socle, and this
socle is contained in the one-dimensional socle of $M$. Consequently
\[
  B \;=\; R/\Ann(H)
\]
is Artinian Gorenstein, has embedding dimension at most three, and has socle
degree $E = e - a$. We write $B_j = \dim_k \bigl( R/\Ann(H) \bigr)_j$ for its
Hilbert function. For every $j \le n$ and $c \in R_j$,
\[
  cH = 0
  \iff cv \in \ker(R^2 \to M)_{a+j}
  \overset{\eqref{eq:seven}}{\iff} cv \in (vJ)_{a+j}
  \iff c \in J_j .
\]
The last equivalence uses the injectivity of multiplication by the nonzero
vector $v$. Thus $J$ and $\Ann(H)$ agree through degree $n$.

\medskip
\noindent\textbf{Equation (8).}
In both the $H = 0$ and $H \ne 0$ cases, with negative subscripts
interpreted as zero, we claim that for $t = d-2, d-1, d$,
\begin{equation} \label{eq:eight}
  g_t \;=\; 2N_t - N_{t-a} + B_{t-a} .
\end{equation}
Indeed, $M$ is the cokernel of $\delta_1$, so counting dimensions in degree
$t$ gives
\[
  g_t \;=\; 2N_t - \dim_k \bigl( \ker(R^2 \to M) \bigr)_t .
\]
For $t \le d$ that kernel is $(vJ)_t$ by \eqref{eq:seven}. Multiplication by
$v$ is an injective map $R \to R^2$, homogeneous of degree $a$, since $v \ne
0$ and $R$ is a domain; hence $(vJ)_t \cong J_{t-a}$ and
$\dim_k (vJ)_t = \dim_k J_{t-a}$. Finally
$\dim_k J_{t-a} = N_{t-a} - \dim_k (R/J)_{t-a}$, and for $t \le d$ we have
$t - a \le n$, so $\dim_k (R/J)_{t-a} = B_{t-a}$ because $J$ and $\Ann(H)$
agree through degree $n$. Combining these gives \eqref{eq:eight}.

\medskip
\noindent\textbf{Conclusion.}
Equation \eqref{eq:eight} gives $g_{d-1} \ge N_{d-1}$, since
$B_{d-1-a} \ge 0$ and $N_{d-1-a} \le N_{d-1}$. On the other hand,
\[
  g_{d-1} \;=\; h_{e-d+1} \;\le\; N_{e-d+1} .
\]
Since $N_j$ is strictly increasing, it follows that
\begin{equation} \label{eq:nine}
  2d \;\le\; e + 2 .
\end{equation}
When $H \ne 0$, this gives
\[
  2(n-1) \;=\; 2d - 2a - 2 \;\le\; e - 2a \;\le\; e - a \;=\; E,
\]
so Lemma \ref{lem:stanley} applies to $B$ at the indices $n-2 \le n-1$ and
gives
\[
  D \;:=\; B_{n-1} - B_{n-2} \;\ge\; 0 .
\]
(If $n \le 1$ then $B_{n-2} = 0$ by the convention on negative subscripts,
and $D \ge 0$ holds because $B_{n-1} \ge 0$, without appeal to Lemma
\ref{lem:stanley}.) For $H = 0$ the inequality $D \ge 0$ also holds, since
$D = 0$ by definition.

Set $x = g_{d-1} - g_{d-2}$. Using \eqref{eq:eight} and
$N_j - N_{j-1} = j+1$, we obtain
\[
  x \;=\; 2d - n + D \;=\; d + a + D \;\ge\; d .
\]
Also $g_{d-2} \le 2N_{d-2} = d(d-1)$.

Finally, \eqref{eq:six} says that the tested triple has the form
\[
  (g_{d-2}, g_{d-1}, g_d) \;=\; (y,\; y+x,\; y+2x+1)
\]
for $y = g_{d-2}$. Its log-concavity margin is
\[
  g_{d-1}^2 - g_{d-2} g_d
  \;=\; (y+x)^2 - y(y + 2x + 1)
  \;=\; x^2 - y
  \;\ge\; d^2 - d(d-1) \;=\; d \;>\; 0 .
\]
This contradicts the assumed failure.

We have proved log-concavity of $g$ at every index $d-1$ with
$2 \le d \le e$. Reversal preserves the log-concavity inequalities, and
$g_d = h_{e-d}$, so the same inequalities hold for $h$.
\end{proof}

\begin{remark}
The argument is valid over every field. The substantive theorem invoked is
the characteristic-free form of Stanley's theorem in Lemma
\ref{lem:stanley}; the remaining steps use graded Matlis duality, linear
algebra over $\Frac(R)$, and the UFD property of $k[x_1,x_2,x_3]$.
\end{remark}

\subsection*{Acknowledgment}

Artificial intelligence was used in producing both the proof and its Lean
formalization.

\begin{thebibliography}{9}

\bibitem{AlphaProofNexus}
Google DeepMind,
\emph{Advancing Mathematics Research with AI-Driven Formal Proof Search},
preprint (2026), \url{https://arxiv.org/abs/2605.22763}.

\bibitem{Iarrobino}
A.\ Iarrobino,
\emph{Log-concave Gorenstein sequences},
preprint, \url{https://arxiv.org/abs/2207.10753}.

\bibitem{Stanley}
R.\ P.\ Stanley,
\emph{Hilbert functions of graded algebras},
Adv.\ Math.\ \textbf{28} (1978), 57--83.

\bibitem{ZanelloStanley}
F.\ Zanello,
\emph{Stanley's theorem on codimension 3 Gorenstein $h$-vectors},
Proc.\ Amer.\ Math.\ Soc.\ \textbf{134} (2006), 5--8,
\url{https://arxiv.org/abs/math/0411231}.

\bibitem{ZanelloLC}
F.\ Zanello,
\emph{Log-concavity of level Hilbert functions and pure $O$-sequences},
J.\ Commut.\ Algebra \textbf{16} (2024), no.\ 2, 245--256,
\url{https://arxiv.org/abs/2210.09447}.

\end{thebibliography}

\vspace{1em}
\noindent\emph{2020 Mathematics Subject Classification.}
Primary 13D40; Secondary 13H10, 13D02, 05E40.

\medskip
\noindent\emph{Key words and phrases.}
Level algebra, Hilbert function, log-concavity, minimal free resolution,
graded Matlis duality, Gorenstein algebra.

\end{document}
